\setcounter{section}{5}

\Needspace{5\baselineskip}
\hypertarget{ux591aux4e13ux5bb6on-policy-distillation}{%
\section{多专家On-policy Distillation}\label{ux591aux4e13ux5bb6on-policy-distillation}}

DeepSeek-V4采用了多专家On-policy Distillation的方法，先基于预训练模型，用SFT和RL训练多个领域专家，再将其知识蒸馏到统一的学生模型中。

\Needspace{5\baselineskip}
\hypertarget{ux4e00ux9886ux57dfux4e13ux5bb6ux6a21ux578bux8badux7ec3}{%
\subsection{一、领域专家模型训练}\label{ux4e00ux9886ux57dfux4e13ux5bb6ux6a21ux578bux8badux7ec3}}

基于预训练模型，用SFT和GRPO训练一批领域专家模型，领域包括数学、代码、Agent、指令跟随等。业界一般还会给每个领域训练多种推理强度的子版本，分别对应Non-think、Think、Think Max，三种模式在RL时用不同的length penalty和context window。

做Agent类专家时还引入了~Quick Instruction~机制，聊天产品里有很多附加任务，比如判断是否触发搜索、识别意图。传统做法是另一个小模型做这些，导致每次上下文变化，要重新进行预填充。V4 的做法是给输入序列直接附一组special token，每个token对应一个附加任务，复用现成的KV cache，降低首字延迟。

\Needspace{5\baselineskip}
\hypertarget{ux4e8con-policy-distillation}{%
\subsection{二、On-Policy Distillation}\label{ux4e8con-policy-distillation}}

运用On-Policy蒸馏方法，把所有专家的知识融合到统一的学生模型里。相比于Off-policy方法而言，On-policy更能避免灾难性遗忘。具体而言，即让学生在自己生成的轨迹上学习多个专家模型的输出概率分布。

\FloatBarrier
\Needspace{5\baselineskip}
\hypertarget{ux53c2ux8003ux6587ux732e-106}{%
\subsection{参考文献}\label{ux53c2ux8003ux6587ux732e-106}}

\vspace{0.25em}
\begin{itemize}
\tightlist
\item
  DeepSeek-AI. (2026). \href{https://arxiv.org/abs/2606.19348}{DeepSeek-V4: Towards Highly Efficient Million-Token Context Intelligence}. arXiv:2606.19348.
\end{itemize}

\clearpage
